% SPDX-License-Identifier: GPL-3.0-or-later
% Copyright (C) 2025-2026 西濱大将 (NISHIHAMA Daisuke)
% Euler vs RANS: near-wall flow and mesh, base treatment, and which axial-force terms come from the solver.
\begin{tikzpicture}[
  font=\footnotesize,
  ttl/.style={font=\small\bfseries, anchor=west},
  lbl/.style={font=\scriptsize, cink2, align=center},
  cfdbox/.style={draw=cblue, fill=cfill, thick, rounded corners=2pt, minimum height=7mm, align=center, font=\scriptsize},
  modbox/.style={draw=corange, fill=cfillo, thick, rounded corners=2pt, minimum height=7mm, align=center, font=\scriptsize},
  vel/.style={->, thick, cblue}]
  % ---------------- column titles
  \node[ttl] at (-0.2,6.4) {(a) Euler（\texttt{model = "euler"}、既定）};
  \node[ttl] at (8.3,6.4) {(b) RANS（\texttt{model = "rans"}）};
  % ---------------- row 1: near-wall flow and mesh
  \node[font=\scriptsize\bfseries, cink2, anchor=west] at (-0.2,5.8) {壁の近くの流れと格子};
  % Euler: tetrahedra up to the wall, uniform velocity (slip)
  \begin{scope}[shift={(0,2.6)}]
    \fill[cfillg] (0,0) rectangle (6.6,-0.35);
    \draw[very thick, cink] (0,0) -- (6.6,0);
    \foreach \x in {0,1.1,...,5.5} {
      \draw[cgrid] (\x,0) -- (\x+0.55,0.95) -- (\x+1.1,0) ;
      \draw[cgrid] (\x+0.55,0.95) -- (\x+1.65,0.95);
      \draw[cgrid] (\x+0.55,0.95) -- (\x+1.1,1.9) -- (\x+1.65,0.95);
      \draw[cgrid] (\x+1.1,1.9) -- (\x+2.2,1.9);
    }
    \foreach \y in {0.12,0.6,1.1,1.6,2.1,2.6} \draw[vel] (0.3,\y) -- (1.5,\y);
    \node[lbl, anchor=west] at (1.7,2.6) {壁まで一様な速度（滑り壁）};
    \node[lbl, anchor=north] at (3.3,-0.4) {四面体だけ。壁では流れが壁を通り抜けないだけで、\\壁に沿う速度は 0 にならない。粘性も摩擦もない};
  \end{scope}
  % RANS: prism layers growing from the wall, boundary-layer profile
  \begin{scope}[shift={(8.5,2.6)}]
    \fill[cfillg] (0,0) rectangle (6.6,-0.35);
    \draw[very thick, cink] (0,0) -- (6.6,0);
    \foreach \h in {0.04,0.09,0.15,0.23,0.33,0.45,0.60,0.78,1.0} \draw[corange!70] (0,\h) -- (6.6,\h);
    \foreach \x in {0.6,1.2,...,6.0} \draw[corange!70] (\x,0) -- (\x,1.0);
    \foreach \x in {0,1.1,...,5.5} {
      \draw[cgrid] (\x,1.0) -- (\x+0.55,1.9) -- (\x+1.1,1.0);
      \draw[cgrid] (\x+0.55,1.9) -- (\x+1.65,1.9);
    }
    \draw[thick, cblue] (0.3,0) .. controls (0.75,0.06) and (1.25,0.35) .. (1.45,1.0) -- (1.5,2.6);
    \foreach \y/\l in {0.1/0.45,0.3/0.85,0.6/1.08,1.0/1.17,1.6/1.2,2.1/1.2,2.6/1.2} \draw[vel] (0.3,\y) -- (0.3+\l,\y);
    \draw[decorate, decoration={brace, amplitude=3pt, mirror}, cink2] (6.75,0) -- (6.75,1.0) node[midway, right=3pt, lbl, anchor=west] {プリズム層\\（境界層）};
    \node[lbl, anchor=west] at (1.7,2.6) {壁で速度 0（粘着壁）、摩擦が生じる};
    \node[lbl, anchor=north] at (3.3,-0.4) {壁の上にプリズム層を積み、その外を四面体で埋める。\\最初の層の高さを \texttt{yplus} で決め、乱流は SA モデルで表す};
  \end{scope}
  % ---------------- row 2: base
  \node[font=\scriptsize\bfseries, cink2, anchor=west] at (-0.2,0.9) {機体の底面};
  \begin{scope}[shift={(0.6,-0.9)}]
    \fill[cfillg, draw=cink, thick] (0,-0.55) -- (3.6,-0.55) -- (3.6,0.55) -- (0,0.55) -- cycle;
    \draw[thick, dashed, cblue, fill=cfill] (3.6,0.55) -- (5.6,0.06) -- (5.6,-0.06) -- (3.6,-0.55);
    \foreach \y in {0.95,-0.95} \draw[vel] (0.2,\y) -- (5.6,\y);
    \node[lbl, anchor=north] at (2.9,-1.15) {底面の後ろに円すいの尾部フェアリングを置く。\\フェアリングの面の力は積分しない};
  \end{scope}
  \begin{scope}[shift={(9.1,-0.9)}]
    \fill[cfillg, draw=cink, thick] (0,-0.55) -- (3.6,-0.55) -- (3.6,0.55) -- (0,0.55) -- cycle;
    \foreach \y in {0.95,-0.95} \draw[vel] (0.2,\y) -- (5.6,\y);
    \draw[thick, corange, ->] (4.15,0.08) arc (-90:240:0.2);
    \draw[thick, corange, ->] (4.15,-0.08) arc (90:-240:0.2);
    \draw[thick, dotted, corange] (3.6,0.55) .. controls (4.8,0.5) .. (5.6,0.25);
    \draw[thick, dotted, corange] (3.6,-0.55) .. controls (4.8,-0.5) .. (5.6,-0.25);
    \node[lbl, anchor=north] at (2.9,-1.15) {平らな底面のまま解き、後ろに剝離した後流ができる。\\底面の力を使うかは \texttt{rans\_base} で選ぶ};
  \end{scope}
  % ---------------- row 3: axial force terms
  \node[font=\scriptsize\bfseries, cink2, anchor=west] at (-0.2,-3.0) {軸力の項の出どころ};
  \begin{scope}[shift={(0,-3.9)}]
    \node[cfdbox, minimum width=27mm, anchor=west] (e1) at (0,0) {圧力（造波抗力を含む）};
    \node[modbox, minimum width=37mm, anchor=west] (e2) at (2.9,0) {摩擦、底面、フィン縁、\\ボートテール、追加抗力};
  \end{scope}
  \begin{scope}[shift={(8.5,-3.9)}]
    \node[cfdbox, minimum width=25mm, anchor=west] (r1) at (0,0) {圧力と摩擦\\（フィン縁を含む）};
    \node[modbox, minimum width=19mm, anchor=west] (r2) at (2.7,0) {底面\\（既定）};
    \node[modbox, minimum width=19mm, anchor=west] (r3) at (4.8,0) {粗さと縁の\\補正、追加抗力};
  \end{scope}
  \node[cfdbox, minimum width=6mm, minimum height=4mm] at (3.2,-5.0) {};
  \node[lbl, anchor=west] at (3.5,-5.0) {CFD の解から得る};
  \node[modbox, minimum width=6mm, minimum height=4mm] at (7.6,-5.0) {};
  \node[lbl, anchor=west] at (8.3,-5.0) {部品積み上げ法の式で補う};
  % separator
  \draw[cgrid, thick] (7.75,6.2) -- (7.75,-4.5);
\end{tikzpicture}
